%%%%%%%%%%%%%%%%%%%%%%%%%%%%%%%%%%%%%%%%%%%%%%%%%%%%%%%%%%%%%%%%%%%%%%%%%%%%%%%%%%
%%% Paper Template — Dual Format (Article / LLNCS)
%%%%%%%%%%%%%%%%%%%%%%%%%%%%%%%%%%%%%%%%%%%%%%%%%%%%%%%%%%%%%%%%%%%%%%%%%%%%%%%%%%

%%% Supported modes:
%%%   • Submission mode  → LLNCS conference version
%%%   • Eprint mode      → full version / arXiv
%%%   • Comments mode    → enables author annotations

%%%%%%%%%%%%%%%%%%%%%%%%%%%%%%%%%%%%%%%%%%%%%%%%%%%%%%%%%%%%%%%%%%%%%%%%%%%%%%%%%%
%%% Build Toggles
%%%%%%%%%%%%%%%%%%%%%%%%%%%%%%%%%%%%%%%%%%%%%%%%%%%%%%%%%%%%%%%%%%%%%%%%%%%%%%%%%%

\newif\ifSubmission
\newif\ifEprint
\newif\ifComments

\Submissiontrue
\Eprintfalse
\Commentsfalse

%%%%%%%%%%%%%%%%%%%%%%%%%%%%%%%%%%%%%%%%%%%%%%%%%%%%%%%%%%%%%%%%%%%%%%%%%%%%%%%%%%
%%% Document Class
%%%%%%%%%%%%%%%%%%%%%%%%%%%%%%%%%%%%%%%%%%%%%%%%%%%%%%%%%%%%%%%%%%%%%%%%%%%%%%%%%%

\ifSubmission
  \documentclass[
    letterpaper,
    runningheads,
    USenglish
  ]{llncs}
\else
  \documentclass[
    11pt,
    letterpaper,
    USenglish
  ]{article}
\fi

%%%%%%%%%%%%%%%%%%%%%%%%%%%%%%%%%%%%%%%%%%%%%%%%%%%%%%%%%%%%%%%%%%%%%%%%%%%%%%%%%%
%%% Package Configuration — Submission Mode
%%%%%%%%%%%%%%%%%%%%%%%%%%%%%%%%%%%%%%%%%%%%%%%%%%%%%%%%%%%%%%%%%%%%%%%%%%%%%%%%%%

\ifSubmission

  \PassOptionsToPackage{inline}{enumitem}

  % Core mathematics and crypto infrastructure
  \usepackage[crypto,captions,tikz,theorems]{llncscrypto}

  % Complexity classes and notation
  \usepackage[full]{complexity}

  % Utility packages
  \usepackage{comment}
  \usepackage[bottom]{footmisc}

  % Tables
  \usepackage{
    array,
    colortbl,
    makecell,
    multirow,
    tabularray,
    tabularx
  }

%%%%%%%%%%%%%%%%%%%%%%%%%%%%%%%%%%%%%%%%%%%%%%%%%%%%%%%%%%%%%%%%%%%%%%%%%%%%%%%%%%
%%% Package Configuration — Article / Eprint Mode
%%%%%%%%%%%%%%%%%%%%%%%%%%%%%%%%%%%%%%%%%%%%%%%%%%%%%%%%%%%%%%%%%%%%%%%%%%%%%%%%%%

\else

  % Core mathematics setup
  \ifEprint
    \usepackage[latinmodern,draft]{csamsmath}
  \else
    \usepackage[latinmodern]{csamsmath}
  \fi

  % Typography
  \usepackage{sansmath}
  \usepackage{setspace}

  % Complexity classes and crypto notation
  \usepackage[full]{complexity}
  \usepackage{tcscrypto}

  % Metadata and references
  \usepackage{orcidlink}
  \usepackage{tocbibind}
  \usepackage{appendix}

  % Utility packages
  \usepackage[bottom]{footmisc}

\fi

%%%%%%%%%%%%%%%%%%%%%%%%%%%%%%%%%%%%%%%%%%%%%%%%%%%%%%%%%%%%%%%%%%%%%%%%%%%%%%%%%%
%%% Shared Packages and Configuration
%%%%%%%%%%%%%%%%%%%%%%%%%%%%%%%%%%%%%%%%%%%%%%%%%%%%%%%%%%%%%%%%%%%%%%%%%%%%%%%%%%

%%% Draft spacing
\ifEprint
  \linespread{1.10}
\fi

%%% Typography helpers
\usepackage[normalem]{ulem}
\usepackage{pifont}

%%%%%%%%%%%%%%%%%%%%%%%%%%%%%%%%%%%%%%%%%%%%%%%%%%%%%%%%%%%%%%%%%%%%%%%%%%%%%%%%%%
%%% Graphics and Diagrams
%%%%%%%%%%%%%%%%%%%%%%%%%%%%%%%%%%%%%%%%%%%%%%%%%%%%%%%%%%%%%%%%%%%%%%%%%%%%%%%%%%

\ifSubmission\else
  \usepackage{amssymb}
  \usepackage{tikz}
  \usepackage{tikz-cd}
  \usepackage{tikz-3dplot}
  \usepackage{pgfplots}

  \pgfplotsset{compat=1.18}
\fi

\usetikzlibrary{
  arrows.meta,
  backgrounds,
  calc,
  chains,
  decorations.markings,
  decorations.pathreplacing,
  fit,
  graphs,
  matrix,
  positioning,
  shapes.geometric,
  shapes.misc
}

%%%%%%%%%%%%%%%%%%%%%%%%%%%%%%%%%%%%%%%%%%%%%%%%%%%%%%%%%%%%%%%%%%%%%%%%%%%%%%%%%%
%%% Compatibility Helpers
%%%%%%%%%%%%%%%%%%%%%%%%%%%%%%%%%%%%%%%%%%%%%%%%%%%%%%%%%%%%%%%%%%%%%%%%%%%%%%%%%%

\makeatletter
\let\cc\relax
\@ifundefined{cc}{\newcommand{\cc}{}}{}
\makeatother

%%%%%%%%%%%%%%%%%%%%%%%%%%%%%%%%%%%%%%%%%%%%%%%%%%%%%%%%%%%%%%%%%%%%%%%%%%%%%%%%%%
%%% Author Comments
%%%%%%%%%%%%%%%%%%%%%%%%%%%%%%%%%%%%%%%%%%%%%%%%%%%%%%%%%%%%%%%%%%%%%%%%%%%%%%%%%%

\ifComments
  \renewcommand{\missing}[1]{%
    \textcolor{WildStrawberry}{%
      \ifSubmission
        {\sffamily #1}
      \else
        {\sffamily\sansmath #1}
      \fi
    }%
    \marginnote{%
      \textcolor{WildStrawberry}{\LARGE$\star$}%
    }%
  }
  % Example:
  % \newcomment{Alice}{RoyalBlue}{alice}
  % \newcomment{Bob}{Firebrick4}{bob}
\else
  \renewcommand{\missing}[1]{#1}
\fi

%%%%%%%%%%%%%%%%%%%%%%%%%%%%%%%%%%%%%%%%%%%%%%%%%%%%%%%%%%%%%%%%%%%%%%%%%%%%%%%%%%
%%% Project Macros
%%%%%%%%%%%%%%%%%%%%%%%%%%%%%%%%%%%%%%%%%%%%%%%%%%%%%%%%%%%%%%%%%%%%%%%%%%%%%%%%%%

% %%%%%%%%%%%%%%%%%%%%%%%%%%%%%%%%%%%%%%%%%%%%%%%%%%%%%%%%%%%%%%%%%%%%%%%%%%%%%%%%%%
%  macros.tex — paper-specific macros
%
%  tcscrypto.sty (bundled) already defines the following — do NOT redefine:
%    Math operators : \Pr, \E, \poly, \polylog, \negl, \log, \lg
%    Sampling       : \getsr, \longsample, \sample
%    Number sets    : \NN, \ZZ, \QQ, \RR, \FF, \GG, \CC
%    Blackboard bold: \bbA–\bbZ (single-letter \bbX shorthand)
%    Calligraphic   : \calA–\calZ (\calX shorthand)
%    Mathsf letters : \sfA–\sfZ (\sfX shorthand)
%    Security param : \secparam, \secpar, \SecParam, \comparam
%    Primitives     : \Gen, \Eval, \Vrfy, \Sign, \Prove, \Verify,
%                     \Enc, \Dec, \Setup, \KeyGen, \Hash, \Sim, \Ext
%    Crypto schemes : \PRF, \PRG, \OWF, \PKE, \SKE, \SIG, \NIZK,
%                     \SNARG, \SNARK, \IOP, \MPC, \FHE, \IBE, ... (100+)
%    Scheme methods : \PRFEval, \PKEEnc, \SIGSign, ... (scheme.Method form)
%    Adversaries    : \Adv, \advA, \advB, \advC, \advD, \advZ, \Red
%    Proof entities : \Prover, \Verifier, \MaliciousProver, ...
%    Security games : \advantage, \newsequenceofgames, \nextgame, \gameref
%    Oracles        : \HashOracle, \SignOracle, \CDHOracle, ...
%    Security notions: \indcpa, \indcca, \eufcma, \sufcma, ...
%    Problems       : \LWE, \SIS, \CDH, \DDH, \RSA, \SATProb, ... (60+)
%    Worlds         : \Minicrypt, \Cryptomania, \Obfustopia, ...
%    Protocol env   : protocolbox
%    Pseudocode     : \pcif, \pcelse, \pcfor, \pcreturn, \pclet, ...
%    Abbreviations  : \eg, \ie, \cf, \wrt, \etal, \vs, \nb
%    Delimiters     : \abs, \floor, \ceil, \norm, \inner, \set, \tuple
%    Distributions  : \Distribution, \Distr, \Ensemble, \cindist, \sindist
%    Lattice        : \latbasis, \latvec, \laterror, \latpubmat, ...
%    Pairing groups : \GGone, \GGtwo, \GGt, \GT, \pairing
%    Kolmogorov     : \Kt, \KS, \KC, \pKt, \rKS, ...
%    Proof complexity: \Axioms, \Clauses, \proves, \Tseitin, ...
%    TFNP classes   : \CLS, \EOL, \EOML, \PWPP, ...
%
%  See tcscrypto.sty for the full reference.
%%%%%%%%%%%%%%%%%%%%%%%%%%%%%%%%%%%%%%%%%%%%%%%%%%%%%%%%%%%%%%%%%%%%%%%%%%%%%%%%%%

%%%%%%%%%%%%%%%%%%%%%%%%%%%%%%%%%%%%%%%%%%%%%%%%%%%%%%%%%%%%%%%%%%%%%%%%%%%%%%%%%%
%  Email (LLNCS cls already defines \email; only define in article mode)
%%%%%%%%%%%%%%%%%%%%%%%%%%%%%%%%%%%%%%%%%%%%%%%%%%%%%%%%%%%%%%%%%%%%%%%%%%%%%%%%%%
\ifSubmission\else
  \newcommand{\email}[1]{\texttt{#1}}
\fi

%%%%%%%%%%%%%%%%%%%%%%%%%%%%%%%%%%%%%%%%%%%%%%%%%%%%%%%%%%%%%%%%%%%%%%%%%%%%%%%%%%
%  Add paper-specific commands below this line
%%%%%%%%%%%%%%%%%%%%%%%%%%%%%%%%%%%%%%%%%%%%%%%%%%%%%%%%%%%%%%%%%%%%%%%%%%%%%%%%%%


%%%%%%%%%%%%%%%%%%%%%%%%%%%%%%%%%%%%%%%%%%%%%%%%%%%%%%%%%%%%%%%%%%%%%%%%%%%%%%%%%%
%%% Document
%%%%%%%%%%%%%%%%%%%%%%%%%%%%%%%%%%%%%%%%%%%%%%%%%%%%%%%%%%%%%%%%%%%%%%%%%%%%%%%%%%

\begin{document}

%%%%%%%%%%%%%%%%%%%%%%%%%%%%%%%%%%%%%%%%%%%%%%%%%%%%%%%%%%%%%%%%%%%%%%%%%%%%%%%%%%
%%% Title and Author Metadata — Submission Mode
%%%%%%%%%%%%%%%%%%%%%%%%%%%%%%%%%%%%%%%%%%%%%%%%%%%%%%%%%%%%%%%%%%%%%%%%%%%%%%%%%%

\ifSubmission
	\title{Paper Title}

	\author{
		First Author\inst{1}
		\and
		Second Author\inst{2}
	}

	\institute{
		Institution One, City, Country\\
		\email{author1@example.com}
		\and
		Institution Two, City, Country\\
		\email{author2@example.com}
	}

	\maketitle

	%%%%%%%%%%%%%%%%%%%%%%%%%%%%%%%%%%%%%%%%%%%%%%%%%%%%%%%%%%%%%%%%%%%%%%%%%%%%%%%%%%
	%%% Title and Author Metadata — Article / Eprint Mode
	%%%%%%%%%%%%%%%%%%%%%%%%%%%%%%%%%%%%%%%%%%%%%%%%%%%%%%%%%%%%%%%%%%%%%%%%%%%%%%%%%%

\else

	\pagestyle{plain}
	\setcounter{tocdepth}{2}

	\ifEprint
		\title{Paper Title (Full Version)}
	\else
		\title{Paper Title}
	\fi

	\ifEprint
		\author[1]{First Author}
		\author[2]{Second Author}

		\affil[1]{
			Institution One, City, Country
			\par\email{author1@example.com}\medskip
		}
		\affil[2]{
			Institution Two, City, Country
			\par\email{author2@example.com}
		}

	\else

		\author{}

	\fi

	\date{\vspace{-2.5em}}

	\maketitle

\fi

%%%%%%%%%%%%%%%%%%%%%%%%%%%%%%%%%%%%%%%%%%%%%%%%%%%%%%%%%%%%%%%%%%%%%%%%%%%%%%%%%%
%%% Abstract
%%%%%%%%%%%%%%%%%%%%%%%%%%%%%%%%%%%%%%%%%%%%%%%%%%%%%%%%%%%%%%%%%%%%%%%%%%%%%%%%%%

\begin{abstract}
	The abstract should communicate the core contribution of the paper in a
	self-contained form. A strong abstract typically answers four questions:
	\begin{itemize}
		\item What problem does the paper study?
		\item Why is the problem important or technically interesting?
		\item What are the main results or techniques?
		\item What limitations, assumptions, or implications matter most?
	\end{itemize}
	State the main theorem or contribution explicitly whenever possible.
	Avoid excessive notation, citations, undefined terminology, and long
	historical background. The abstract should remain readable to researchers
	outside the immediate specialization.
\end{abstract}

%%%%%%%%%%%%%%%%%%%%%%%%%%%%%%%%%%%%%%%%%%%%%%%%%%%%%%%%%%%%%%%%%%%%%%%%%%%%%%%%%%
%%% Introduction
%%%%%%%%%%%%%%%%%%%%%%%%%%%%%%%%%%%%%%%%%%%%%%%%%%%%%%%%%%%%%%%%%%%%%%%%%%%%%%%%%%

\section{Introduction}
\label{sec:introduction}

The introduction should move from a broad context to the precise technical
problem addressed in the paper.

Explain:
\begin{itemize}
	\item the scientific or mathematical setting,
	\item the motivating question,
	\item the main technical obstacles,
	\item the significance of resolving them.
\end{itemize}

A good introduction helps readers understand not only \emph{what} the
paper proves, but also \emph{why} the result matters and how it fits
into the existing literature.

Clearly distinguish:
\begin{itemize}
	\item established prior results,
	\item open problems or limitations,
	\item the new ideas introduced in this work.
\end{itemize}

Avoid delaying the main contribution for too long. Readers should
understand the central theorem and its relevance early in the section.

%%%%%%%%%%%%%%%%%%%%%%%%%%%%%%%%%%%%%%%%%%%%%%%%%%%%%%%%%%%%%%%%%%%%%%%%%%%%%%%%%%
%%% Contributions
%%%%%%%%%%%%%%%%%%%%%%%%%%%%%%%%%%%%%%%%%%%%%%%%%%%%%%%%%%%%%%%%%%%%%%%%%%%%%%%%%%

\subsection{Our Contributions}
\label{sec:contributions}

This subsection should summarize the paper's contributions in a concise
and technically precise form.

Each item should state:
\begin{itemize}
	\item the result,
	\item the assumptions under which it holds,
	\item the main novelty or improvement.
\end{itemize}

Organize contributions from most central to most technical. Avoid vague
claims such as ``significant improvement'' unless quantified.

\begin{itemize}
	\item \textbf{Main theorem.}
	      State the primary result and the setting in which it applies.
	\item \textbf{Technical contribution.}
	      Explain the key technique, reduction, construction, or conceptual
	      innovation introduced by the paper.
\end{itemize}

When appropriate, briefly compare quantitative parameters with prior
work.

%%%%%%%%%%%%%%%%%%%%%%%%%%%%%%%%%%%%%%%%%%%%%%%%%%%%%%%%%%%%%%%%%%%%%%%%%%%%%%%%%%
%%% Related Work
%%%%%%%%%%%%%%%%%%%%%%%%%%%%%%%%%%%%%%%%%%%%%%%%%%%%%%%%%%%%%%%%%%%%%%%%%%%%%%%%%%

\subsection{Related Work}
\label{sec:related}

The related work section should help readers position the paper within
the broader literature.

Focus on:
\begin{itemize}
	\item papers most directly connected to the main results,
	\item differences in assumptions, models, or guarantees,
	\item conceptual distinctions between techniques.
\end{itemize}

Avoid turning the section into a chronological list of references. Prioritize
clarity and relevance over exhaustive citation lists.

When comparing with previous work, explain precisely:
\begin{itemize}
	\item what earlier papers achieved,
	\item what remained unresolved,
	\item what this work changes or improves.
\end{itemize}

%%%%%%%%%%%%%%%%%%%%%%%%%%%%%%%%%%%%%%%%%%%%%%%%%%%%%%%%%%%%%%%%%%%%%%%%%%%%%%%%%%
%%% Preliminaries
%%%%%%%%%%%%%%%%%%%%%%%%%%%%%%%%%%%%%%%%%%%%%%%%%%%%%%%%%%%%%%%%%%%%%%%%%%%%%%%%%%

\section{Preliminaries}
\label{sec:prelim}

Introduce notation, conventions, and background material required for
the technical sections.

Definitions should:
\begin{itemize}
	\item use consistent notation,
	\item avoid unnecessary generality,
	\item isolate assumptions clearly.
\end{itemize}

Whenever possible, separate standard material from paper-specific
definitions so readers can distinguish foundational concepts from new
technical machinery.

\begin{definition}[Example Definition]
	\label{def:example}
	A \emph{widget} is a tuple $(\Gen,\Eval,\Vrfy)$ of PPT algorithms
	satisfying the following properties \ldots
\end{definition}

%%%%%%%%%%%%%%%%%%%%%%%%%%%%%%%%%%%%%%%%%%%%%%%%%%%%%%%%%%%%%%%%%%%%%%%%%%%%%%%%%%
%%% Main Results
%%%%%%%%%%%%%%%%%%%%%%%%%%%%%%%%%%%%%%%%%%%%%%%%%%%%%%%%%%%%%%%%%%%%%%%%%%%%%%%%%%

\section{Main Results}
\label{sec:main}

Present the primary technical statements early and clearly.

The beginning of the section should help readers understand:
\begin{itemize}
	\item the exact theorem statements,
	\item the proof strategy at a high level,
	\item the relationship between intermediate lemmas and final results.
\end{itemize}

State assumptions explicitly and avoid hiding important conditions inside
proofs.

\begin{theorem}
	\label{thm:main}
	Under assumption $X$, scheme $Y$ satisfies property $Z$.
\end{theorem}

Before presenting a long proof, briefly explain:
\begin{itemize}
	\item the main reduction or construction,
	\item the source of the key difficulty,
	\item the central technical insight.
\end{itemize}

\begin{proof}
	Construct a reduction $\mathcal{R}$ and analyze its correctness and
	security properties \ldots
\end{proof}

%%%%%%%%%%%%%%%%%%%%%%%%%%%%%%%%%%%%%%%%%%%%%%%%%%%%%%%%%%%%%%%%%%%%%%%%%%%%%%%%%%
%%% Conclusion
%%%%%%%%%%%%%%%%%%%%%%%%%%%%%%%%%%%%%%%%%%%%%%%%%%%%%%%%%%%%%%%%%%%%%%%%%%%%%%%%%%

\section{Conclusion}
\label{sec:conclusion}

The conclusion should summarize the conceptual significance of the work
rather than merely restating theorem statements.

Good conclusions often discuss:
\begin{itemize}
	\item remaining open problems,
	\item limitations of the current techniques,
	\item possible future directions,
	\item broader implications of the results.
\end{itemize}

%%%%%%%%%%%%%%%%%%%%%%%%%%%%%%%%%%%%%%%%%%%%%%%%%%%%%%%%%%%%%%%%%%%%%%%%%%%%%%%%%%
%%% Acknowledgments (Eprint / Full Version Only)
%%%%%%%%%%%%%%%%%%%%%%%%%%%%%%%%%%%%%%%%%%%%%%%%%%%%%%%%%%%%%%%%%%%%%%%%%%%%%%%%%%

\ifEprint
	\vfill
	\clearpage

	\parhead{Acknowledgments.}

	Acknowledge collaborators, discussions, institutional support,
	anonymous reviewers, and funding sources where appropriate.

	Keep acknowledgments concise and professional. Avoid excessively
	informal remarks or long historical narratives.

	\vfill
	\clearpage
\fi

%%%%%%%%%%%%%%%%%%%%%%%%%%%%%%%%%%%%%%%%%%%%%%%%%%%%%%%%%%%%%%%%%%%%%%%%%%%%%%%%%%
%%% Bibliography
%%%%%%%%%%%%%%%%%%%%%%%%%%%%%%%%%%%%%%%%%%%%%%%%%%%%%%%%%%%%%%%%%%%%%%%%%%%%%%%%%%

%%% Recommended bibliography workflow:
%%%
%%% 1. Download CryptoBib:
%%%      https://cryptobib.di.ens.fr/
%%%
%%%    Recommended files:
%%%      • bib/abbrev0.bib
%%%      • bib/crypto.bib
%%%
%%% 2. Download CSTheoryBib:
%%%      https://cstheory.bib.website/
%%%
%%%    Recommended file:
%%%      • bib/cstheory.bib
%%%
%%% 3. Store project-specific references in:
%%%      • bib/bibliography.bib
%%%
%%% 4. Adjust the \bibliography command below as needed.

\ifSubmission
	\bibliographystyle{splncs04}

\else

	\begin{small}
		\vfill
		\clearpage
		\bibliographystyle{alphaurl}
		\fi
		\bibliography{bib/bibliography}
		%%% Full bibliography setup:
		%%% \bibliography{
		%%%     bib/abbrev0,
		%%%     bib/crypto,
		%%%     bib/cstheory,
		%%%     bib/bibliography
		%%% }
		\ifSubmission\else
		\vfill
		\clearpage
	\end{small}

\fi

%%%%%%%%%%%%%%%%%%%%%%%%%%%%%%%%%%%%%%%%%%%%%%%%%%%%%%%%%%%%%%%%%%%%%%%%%%%%%%%%%%
%%% Appendix
%%%%%%%%%%%%%%%%%%%%%%%%%%%%%%%%%%%%%%%%%%%%%%%%%%%%%%%%%%%%%%%%%%%%%%%%%%%%%%%%%%

\appendix

%%%%%%%%%%%%%%%%%%%%%%%%%%%%%%%%%%%%%%%%%%%%%%%%%%%%%%%%%%%%%%%%%%%%%%%%%%%%%%%%%%
%%% Deferred Proofs
%%%%%%%%%%%%%%%%%%%%%%%%%%%%%%%%%%%%%%%%%%%%%%%%%%%%%%%%%%%%%%%%%%%%%%%%%%%%%%%%%%

\section{Deferred Proofs}
\label{appx:proofs}

Place proofs omitted from the main body here.

Appendices should contain:
\begin{itemize}
	\item deferred technical arguments,
	\item auxiliary lemmas,
	\item extended constructions,
	\item additional experimental or computational details when relevant.
\end{itemize}

Whenever possible, maintain the same notation and theorem numbering
conventions used in the main text.

\end{document}